% =====================================================================
% Page 1 - Top page (Annex-1)
% =====================================================================
\begin{center}
    {\fontsize{18}{27}\selectfont\bfseries \PROJECTTITLE\\}
    \vspace{0.6cm}

    {\fontsize{14}{21}\selectfont\bfseries \reporttype\\}
    \vspace{0.7cm}

    {\fontsize{14}{21}\selectfont\itshape Submitted by\\}
    \vspace{0.3cm}

    {\fontsize{14}{21}\selectfont\bfseries \studentA\\ \studentB\\ \studentC\\}
    \vspace{0.8cm}

    {\fontsize{16}{24}\selectfont\bfseries Of\\ Bachelor of Technology\\}
    \vspace{0.8cm}

    {\fontsize{12}{21}\selectfont\textbf{\textit{Under the guidance of}}\\}
    {\fontsize{14}{21}\selectfont\bfseries \supervisor\\ \supervisordesig\\}
    \vspace{0.8cm}

    \includegraphics[width=0.36\textwidth]{rgukt_logo.jpeg}
    \vspace{0.6cm}

    {\fontsize{14}{16.8}\selectfont\bfseries DEPARTMENT OF COMPUTER SCIENCE AND ENGINEERING\\
    RAJIV GANDHI UNIVERSITY OF KNOWLEDGE TECHNOLOGIES\\
    BASAR, NIRMAL (DIST.),\\
    TELANGANA -- 504107\\}
\end{center}
\clearpage

% =====================================================================
% Page 2 - Cover page (Annex-2)
% =====================================================================
\begin{center}
    {\fontsize{16}{24}\selectfont\bfseries \projecttitle\\}
    \vspace{0.8cm}

    {\fontsize{12}{18}\selectfont\bfseries\itshape Project Report submitted to\\
        Rajiv Gandhi University of Knowledge Technologies, Basar\\
        for the partial fulfillment of the requirements\\
        for the award of the degree of\\}
    \vspace{0.8cm}

    {\fontsize{14}{21}\selectfont\bfseries Bachelor of Technology\\ in\\ Computer Science \& Engineering\\}
    \vspace{0.3cm}

    {\fontsize{12}{18}\selectfont\bfseries\itshape by\\}
    \vspace{0.3cm}

    {\fontsize{14}{21}\selectfont\bfseries \studentA\\ \studentB\\ \studentC\\}
    \vspace{0.5cm}

    {\fontsize{12}{21}\selectfont\bfseries Under the Guidance of\\}
    \vspace{0.3cm}
    {\fontsize{14}{21}\selectfont\bfseries\itshape \supervisor\\ \supervisordesig\\}
    \vspace{0.5cm}

    \includegraphics[width=0.25\textwidth]{rgukt_logo.jpeg}
    \vspace{0.5cm}

    {\fontsize{14}{16.8}\selectfont\bfseries DEPARTMENT OF COMPUTER SCIENCE AND ENGINEERING\\
        RAJIV GANDHI UNIVERSITY OF KNOWLEDGE TECHNOLOGIES, BASAR\\
        \submissionmonth\\}
\end{center}
\clearpage

% =====================================================================
% Page 3 - Certificate (Annex-3)
% =====================================================================
\phantomsection
\addcontentsline{toc}{chapter}{Certificate}
\begin{center}
    \includegraphics[width=0.25\textwidth]{rgukt_logo.jpeg}\\[0.4cm]
    {\fontsize{14}{16.8}\selectfont\bfseries DEPARTMENT OF COMPUTER SCIENCE AND ENGINEERING\\
    RAJIV GANDHI UNIVERSITY OF KNOWLEDGE TECHNOLOGIES, BASAR\\}
    \vspace{1cm}
    {\fontsize{16}{24}\selectfont\bfseries CERTIFICATE\\}
\end{center}

\vspace{0.5cm}

This is to certify that the Project Report entitled \textbf{``\projecttitle''}
submitted by \textbf{\studentA, \studentB} and \textbf{\studentC}, Department of
Computer Science and Engineering, Rajiv Gandhi University of Knowledge Technologies,
Basar, for partial fulfillment of the requirements for the award of Bachelor of
Technology in Computer Science and Engineering, is a bonafide record of the work and
investigations carried out by them under my supervision and guidance.

\vspace{2.5cm}

\begin{minipage}{0.45\textwidth}
    \begin{center}
        \textbf{PROJECT SUPERVISOR}\\
        \vspace{0.5cm}
        \textbf{\supervisor}\\
        \textbf{\supervisordesig}
    \end{center}
\end{minipage}
\hfill
\begin{minipage}{0.45\textwidth}
    \begin{center}
        \textbf{HEAD OF DEPARTMENT}\\
        \vspace{0.5cm}
        \textbf{Dr. B. Venkat Raman}\\
        \textbf{Assistant Professor}
    \end{center}
\end{minipage}

\vspace{1.5cm}
\begin{center}
    \textbf{EXTERNAL EXAMINER}
\end{center}
\clearpage

% =====================================================================
% Page 4 - Declaration (Annex-4)
% =====================================================================
\phantomsection
\addcontentsline{toc}{chapter}{Declaration}
\begin{center}
    \includegraphics[width=0.25\textwidth]{rgukt_logo.jpeg}\\[0.4cm]
    {\fontsize{14}{16.8}\selectfont\bfseries DEPARTMENT OF COMPUTER SCIENCE AND ENGINEERING\\
    RAJIV GANDHI UNIVERSITY OF KNOWLEDGE TECHNOLOGIES, BASAR\\}
    \vspace{1cm}
    {\fontsize{16}{24}\selectfont\bfseries DECLARATION\\}
\end{center}

\vspace{0.8cm}

We hereby declare that the work which is being presented in this project entitled
``\textbf{\projecttitle}'' submitted to \textbf{RAJIV GANDHI UNIVERSITY OF KNOWLEDGE
TECHNOLOGIES, BASAR} in the partial fulfillment of the requirements for the award of
the degree of \textbf{BACHELOR OF TECHNOLOGY} in \textbf{COMPUTER SCIENCE AND
ENGINEERING}, is an authentic record of our own work carried out under the supervision
of \textbf{\supervisor}, \supervisordesig{} in the Department of Computer Science and
Engineering, RGUKT, Basar.

The matter embodied in this project report has not been submitted by us for the award
of any other degree.

\vspace{2cm}
Place: Basar\\
Date :

\begin{flushright}
    \studentA\\
    \studentB\\
    \studentC
\end{flushright}
\clearpage

% =====================================================================
% Acknowledgement
% =====================================================================
\chapter*{Acknowledgement}
\addcontentsline{toc}{chapter}{Acknowledgement}

We would like to express our sincere gratitude to our project supervisor,
\textbf{\supervisor}, \supervisordesig, Department of Computer Science and Engineering,
whose knowledge and guidance motivated us throughout this work. Our supervisor has
consistently been a source of motivation, encouragement and inspiration, and the time
we spent working under this supervision has truly been a pleasure.

We thank our Head of the Department, \textbf{Dr. B. Venkat Raman}, for the support and
guidance provided, and all senior faculty members of the CSE Department for their help
during our course. We also thank the programmers and non-teaching staff of the CSE
Department.

We express our sincere gratitude to our Project Co-ordinator, \textbf{Mr. R. Rahul}, for the guidance and encouragement given during the project review process.

We thank our Vice-Chancellor, \textbf{Prof. A. Govardhan}, and the Management for
providing excellent facilities to carry out our project work.

Finally, special thanks to our parents for their support and encouragement throughout
our life and this course, and to all our friends and well-wishers for their constant
support.

\vspace{1cm}
\begin{flushright}
    \textbf{\studentA}\\
    \textbf{\studentB}\\
    \textbf{\studentC}
\end{flushright}
\clearpage

% =====================================================================
% Abstract
% =====================================================================
\chapter*{Abstract}
\addcontentsline{toc}{chapter}{Abstract}

Containers and Kubernetes have become the standard way to deploy applications, but
most of their security tooling works before deployment (image scanning, admission
policies) and offers little visibility into what a workload actually does while it
runs. An attacker who obtains a shell inside a pod can open reverse shells, steal
service-account tokens, mine crypto-currency or attempt a container escape without
being noticed. This project presents an \textbf{eBPF-based runtime Intrusion Detection
System (IDS) for Kubernetes} that observes kernel activity of every container,
classifies it with machine learning and deterministic rules, and raises explainable
alerts in real time.

An eBPF program written in C with libbpf CO-RE is attached to 16 system-call
tracepoints (process execution, sensitive file opens, network connections,
privilege and namespace changes, module loading, fileless execution). Events are
streamed through a BPF ring buffer to a Python agent deployed as a Kubernetes
DaemonSet, which maps each event to its Pod, container, namespace and workload,
derives 74 behavioural features and scores them with a 10-class Random Forest
(benign and nine attack categories) together with 19 high-precision rules. The result
is a 0--100 risk score with a Low / Medium / High severity and a MITRE ATT\&CK
technique. Alerts and events are stored in MongoDB, served by a FastAPI backend and
visualised in a React dashboard.

The model was trained on a labelled runtime-telemetry dataset of 40,000 events from 52
benign and attack template families with deliberate edge cases, extended with 25,000
real benign kernel events from the public BETH dataset. On a held-out synthetic test
set it reaches an F1-score of 0.992 with a false-positive rate of 0.09\%. Evaluation
on real-world data (BETH honeypot hosts and a recorded Kubernetes container escape)
exposes the generalisation gap of synthetic training data, which is analysed in
detail. In a live minikube cluster the complete system detected 20 out of 20 attack
re-enactments with an average detection latency of 123~ms, while the agent used about
3\% of one CPU core and 200~MB of memory with zero dropped kernel events.

\textbf{Keywords:} eBPF, Kubernetes, container security, intrusion detection, runtime
security, machine learning, Random Forest, MITRE ATT\&CK.
\clearpage

% =====================================================================
% Contents, figures, tables, abbreviations
% =====================================================================
\renewcommand{\contentsname}{Table of Contents}
\tableofcontents
\clearpage

\phantomsection
\addcontentsline{toc}{chapter}{List of Figures}
\listoffigures
\clearpage

\phantomsection
\addcontentsline{toc}{chapter}{List of Tables}
\listoftables
\clearpage

\chapter*{List of Abbreviations and Symbols Used}
\addcontentsline{toc}{chapter}{List of Abbreviations and Symbols Used}

\begin{longtable}{p{3cm}p{10cm}}
    \textbf{API}     & Application Programming Interface \\
    \textbf{ATT\&CK} & Adversarial Tactics, Techniques and Common Knowledge (MITRE) \\
    \textbf{AUC}     & Area Under the (ROC) Curve \\
    \textbf{BPF}     & Berkeley Packet Filter \\
    \textbf{BTF}     & BPF Type Format \\
    \textbf{CNCF}    & Cloud Native Computing Foundation \\
    \textbf{CO-RE}   & Compile Once -- Run Everywhere \\
    \textbf{CPU}     & Central Processing Unit \\
    \textbf{CV}      & Cross-Validation \\
    \textbf{DFD}     & Data Flow Diagram \\
    \textbf{eBPF}    & extended Berkeley Packet Filter \\
    \textbf{FN / FP} & False Negative / False Positive \\
    \textbf{FPR}     & False-Positive Rate \\
    \textbf{HIDS}    & Host-based Intrusion Detection System \\
    \textbf{IDS}     & Intrusion Detection System \\
    \textbf{JSON}    & JavaScript Object Notation \\
    \textbf{JWT}     & JSON Web Token \\
    \textbf{K8s}     & Kubernetes \\
    \textbf{MCC}     & Matthews Correlation Coefficient \\
    \textbf{ML}      & Machine Learning \\
    \textbf{NIDS}    & Network-based Intrusion Detection System \\
    \textbf{PID}     & Process Identifier \\
    \textbf{RBAC}    & Role-Based Access Control \\
    \textbf{REST}    & Representational State Transfer \\
    \textbf{RF}      & Random Forest \\
    \textbf{ROC}     & Receiver Operating Characteristic \\
    \textbf{RSS}     & Resident Set Size (memory) \\
    \textbf{SUID}    & Set-User-ID permission bit \\
    \textbf{TN / TP} & True Negative / True Positive \\
    \textbf{TTL}     & Time To Live \\
    \textbf{UI}      & User Interface \\
    \textbf{UML}     & Unified Modeling Language \\
    $P(\mathrm{mal})$ & Probability that an event is malicious, $1 - P(\mathrm{benign})$ \\
    $\tau$           & Decision threshold on $P(\mathrm{mal})$ (0.5) \\
\end{longtable}
\clearpage
